% 第三章：系统需求分析
% 草稿来源：drafts/final-term/sections/03-system-requirements.md

\section{问题背景}

蛋白质设计任务并非单一模型调用即可完成的孤立计算。一次典型的 de novo 设计过程涉及目标描述、候选序列生成、结构预测、质量筛选、目标评分和结果汇总等多个环节，这些环节之间的数据依赖、工具差异和失败模式使得简单固定流水线难以胜任。具体而言：序列生成工具提供候选空间，结构预测工具给出折叠置信度，质量控制工具负责过滤不可用结果——但不同工具在输入输出格式、运行成本、失败特征和适用边界上差异显著。如果系统仅在任务开始前选定一条工具链并按序执行，一旦中间步骤的结果偏离预期、高代价工具调用失败或质量门禁拒绝大部分候选，系统缺乏有效的应对手段。

本课题面向 de novo 蛋白质设计场景，目标是构建一个能够根据任务约束自动生成工具链、在运行时根据中间证据调整执行路径、并在失败或风险升高时触发可控恢复的智能工作流系统。系统不替代具体的生物学模型，而是为多个蛋白质设计工具之间建立可组合、可执行、可恢复、可审计的工作流控制层。

\section{用户与使用场景}

% 内容同 Markdown 草稿 §3.2

\section{功能需求}

\subsection{任务接入与约束确认}

% 内容同 Markdown 草稿 §3.3.1

\subsection{候选计划生成与工具链选择}

% 内容同 Markdown 草稿 §3.3.2

\subsection{工作流执行与工具适配}

% 内容同 Markdown 草稿 §3.3.3

\subsection{人在环路与决策审查}

% 内容同 Markdown 草稿 §3.3.4

\subsection{运行时自适应与恢复决策}

% 内容同 Markdown 草稿 §3.3.5

\subsection{结果汇总、报告与审计}

% 内容同 Markdown 草稿 §3.3.6

\section{非功能需求}

% 内容同 Markdown 草稿 §3.4

\section{本章小结}

% 内容同 Markdown 草稿 §3.5
